\documentclass[11pt]{article}
\usepackage[T1]{fontenc}
\usepackage{lmodern,amsmath,amssymb,amsthm,mathtools,microtype}
\usepackage[margin=27mm]{geometry}
\usepackage[hidelinks]{hyperref}
\newtheorem{theorem}{Theorem}
\newtheorem{lemma}{Lemma}
\newtheorem{proposition}{Proposition}
\newcommand{\R}{\mathbb R}
\newcommand{\E}{\mathcal E}
\newcommand{\Tr}{\operatorname{Tr}}
\newcommand{\rank}{\operatorname{rank}}
\newcommand{\nct}{\mathsf n}
\title{Critical inverse moments and local observation geometry on finite flat lattices}
\author{Lluis Eriksson}
\date{4 November 2026: consolidated revision}
\begin{document}
\maketitle
\begin{abstract}
This appendix consolidates one sufficient coefficient condition and its limitations in the
existing flat-lattice observation-rank dossier. A global inverse-conductance moment of
order $D/2$ suffices for the mean-budget rank order through an actual real spectral
observation. Every lower positive inverse exponent admits explicit divergence-free
counterexamples. At the critical exponent, spatial observations by means on a full dyadic
cube tree can require logarithmically many blocks even for a fixed weak region; the same
profiles admit a constant-size unrestricted rectangular partition. The finite Fourier
and flow adapters are proved analytically, and exact finite diagnostics are recorded
separately from formal verification. A four-cycle counterexample now has a focal Lean
proof: coercivity on the zero-mean kernel does not justify the same fixed penalty.
The distinction from a freely chosen penalty is essential. The spectral mechanism is classical; independent
novelty, an interacting gauge-theory gap and a continuum limit are not established.
This is a consolidation appendix, not a separate paper or AIRR deposit.
\end{abstract}

\section{Actual objects and the original observation problem}
Fix $D\geq3$, $N_c\geq2$, $c>0$, and the unit-spacing torus
$X_n=(\mathbb Z/n\mathbb Z)^D$. The actual real adjoint coefficient space has dimension
$g=N_c^2-1$; there are $D$ field components, so write $h=Dg$.
For the actual forward difference $d_i u(x)=u(x+e_i)-u(x)$, take one scalar total
coefficient $a(x)=1+w_{i,x,j}>0$ shared by all derivative directions and field components.
Use counting norms and the actual componentwise flat energy
\begin{equation}\label{eq:energy}
 Q_w(U)=\E_a(U)=\sum_{x,i,j}a(x)\|d_iU_j(x)\|^2.
\end{equation}
The existing flat Hodge identity identifies these forms; its formal source is listed
in Section~\ref{sec:evidence}. All new adapters in this appendix remain analytic.
There is no substitution of a covariant or interacting derivative in \eqref{eq:energy}.

The original local data consist of a differentiable map $C$ at zero, with $C(0)=0$,
into an actual finite-dimensional real normed space $F$, and a penalty satisfying
$\Phi(y)\leq K\|y\|^2$ near zero, such that
\begin{equation}\label{eq:local}
 c\|U\|^2\leq Q_w(U)+\Phi(C(U))
\end{equation}
for sufficiently small actual divergence-zero fields. ``Local'' here means a
configuration-space neighbourhood; spatial support of $C$ is an additional condition.
The constant $c$ is fixed, while observation maps and neighbourhoods may depend on $n$.

A useful necessary condition follows directly. If an actual $J$-dimensional
divergence-zero subspace has $Q_w(U)\leq\lambda\|U\|^2$ with $\lambda<c$, then
$\dim F\geq J$. Otherwise a nonzero $U$ in that subspace lies in $\ker DC(0)$.
Along $tU$, $C(tU)=o(t)$ and the penalty is at most $o(t^2)$. Dividing
\eqref{eq:local} by $t^2$ and taking $t\to0$ contradicts $\lambda<c$.
No rank or band certificate is supplied as an extra hypothesis in this argument.

\section{The critical global moment theorem}
Put $p=D/2>1$. For $m\geq1$ suppose that
\begin{equation}\label{eq:moment}
 n^{-D}\sum_{x\in X_n}(m/a(x))^p\leq M<\infty.
\end{equation}
\begin{theorem}\label{thm:critical}
Under \eqref{eq:moment}, there is an actual real linear spectral observation $C$ such that
\[
 c\|U\|^2\leq Q_w(U)+c\|C(U)\|^2
 \quad\hbox{for every actual cochain }U,
\]
and
\begin{equation}\label{eq:rank}
 \dim F\leq C_*\max\{1,(n/\sqrt{1+m})^D\},
\end{equation}
where, with $I_p=(p-1)^2/[p(p+1)]$ and
$\Gamma_p=\frac p{p-1}I_p^{-(p-1)}$, one safe choice is
\[
 C_*=h\left[1+\Gamma_p(3/4)^D D^{p+1}c^pM2^p\right].
\]
Thus $\Phi(y)=c\|y\|^2$, $K=c$ and any positive neighbourhood radii supply the original
local data. In particular the assertion holds on the actual divergence-zero subspace.
\end{theorem}
The original subclass $1\leq a\leq2m$ with exact arithmetic mean $m$ is included,
but those extra normalizations are not needed for sufficiency. The distribution
condition \eqref{eq:moment} remains additional; it is not derived from an arithmetic mean.

\subsection{A finite normalized Fourier lemma}
For a finite abelian group of size $V$, let $\mathcal F$ be its counting-unitary Fourier
matrix, whose entries have modulus $V^{-1/2}$. For a frequency multiplier $g$ put
\[
 L=V^{-1}\sup_{t>0}t^p\#\{k:|g(k)|^2>t\}.
\]
For a nonnegative matrix $B$, write $\nct(s,B)$ for its number of eigenvalues strictly
greater than $s$, counted with ordinary multiplicity.
\begin{lemma}\label{lem:finite}
For $K=M_f\mathcal F^*M_g\mathcal F$ and $s>0$,
\[
 \nct(s,K^*K)\leq\Gamma_p Ls^{-p}\sum_x|f(x)|^{2p}.
\]
The trace and eigenvalue counts are not volume-normalized.
\end{lemma}
\begin{proof}
Remove phases by unitaries and set $A=\mathcal F^*M_{|g|^2}\mathcal F$.
If $L=0$ the assertion is immediate. Otherwise let $B_0=pL/(p-1)$ and let $A^+$
be the inverse on $\operatorname{ran}A$, zero on its kernel. For $0\leq\gamma\leq A$
write $\rho(x)=\langle\delta_x,\gamma\delta_x\rangle$. Positivity makes $\gamma$
annihilate $\ker A$. On the range let $P_E=\mathbf1_{\{A^+\leq E\}}$ and
$R_E=\mathbf1_{\{A^+>E\}}$. The weak-frequency tail gives
\[
 \langle\delta_x,P_EAP_E\delta_x\rangle
 =V^{-1}\sum_{|g(k)|^2\geq1/E}|g(k)|^2\leq B_0E^{p-1}.
\]
Indeed the tail above $t$ is $t$ times its count plus the integral of counts above $t$;
the resulting bound is $VLp\,t^{1-p}/(p-1)$. Non-strict thresholds follow by limits.
The triangle inequality for $\gamma^{1/2}\delta_x$ split by $P_E+R_E$ yields
\[
 \langle\delta_x,R_E\gamma R_E\delta_x\rangle
 \geq\bigl(\sqrt{\rho(x)}-\sqrt{B_0}E^{(p-1)/2}\bigr)_+^2.
\]
Exact spectral layer integration consequently gives
\begin{equation}\label{eq:density}
 \Tr(A^+\gamma)\geq I_pB_0^{-1/(p-1)}\sum_x\rho(x)^{p/(p-1)}.
\end{equation}
The scalar integral reduces to
$\int_0^1(1-t^{(p-1)/2})^2dt=I_p$.
Now take $P$ to be the spectral projection of $A^{1/2}M_{|f|^2}A^{1/2}$ above $s$,
with rank $\nu>0$, and set $\gamma=A^{1/2}PA^{1/2}$. Then
$0\leq\gamma\leq A$, $\Tr(A^+\gamma)\leq\nu$, and
\[
 s\nu<\sum_x|f(x)|^2\rho(x)
 \leq\Bigl(\sum_x|f(x)|^{2p}\Bigr)^{1/p}
       \Bigl(\sum_x\rho(x)^{p/(p-1)}\Bigr)^{(p-1)/p}.
\]
Insert \eqref{eq:density} and raise to $p$. The constant becomes
$B_0I_p^{-(p-1)}=\Gamma_pL$. This also explicitly handles a zero-frequency kernel.
\end{proof}
The density method is classical and follows the Rumin--Frank approach
\cite[Section 2]{frank}. The proof above supplies its finite normalization and kernel
adapter; it is not a claim of a new Cwikel mechanism or independent discovery.

\subsection{Actual lattice frequencies and the weighted inverse}
For $n\geq2$ let $L_0=d^*d$ and work first on the actual counting-mean-zero scalar
space $H_0$. Its Fourier symbols are
\[
 z_i(k)=e^{2\pi\mathrm i k_i/n}-1,\qquad
 \lambda(k)=4\sum_i\sin^2(\pi k_i/n).
\]
Choose centred integer representatives $r_i$ with $|r_i|\leq n/2$, using only one
representative at an even endpoint. Concavity of sine implies
\[
 \lambda(k)\geq16|r|^2/n^2,\qquad
 |g_i(k)|:=|z_i(k)/\lambda(k)|\leq n/(4|r|),\quad k\ne0.
\]
Set $g_i(0)=0$. If $|g_i|>\tau$, then $0<|r|<R=n/(4\tau)$.
For $R\leq1$ there are none, and for $R>1$ there are at most $(3R)^D$ representatives.
Thus
\begin{equation}\label{eq:weak}
 n^{-D}\sup_{t>0}t^{D/2}\#\{|g_i|^2>t\}\leq(3/4)^D.
\end{equation}
This is uniform in $n$, including $n=2$, not an imported infinite-volume bound.

Let $L_a=d^*M_ad$. On $H_0$ it is positive definite. For $f\in H_0$, the flow
$j_*=M_adL_a^{-1}f$ minimizes $\langle j,M_{a^{-1}}j\rangle$ among $d^*j=f$:
write $j=j_*+k$, $d^*k=0$, and the cross term is
$\langle dL_a^{-1}f,k\rangle=0$. Inserting the explicit feasible flow
$j=dL_0^{-1}f$ gives the actual quadratic-form comparison
\begin{equation}\label{eq:inverse}
 L_a^{-1}\leq\sum_{i=1}^DT_i^*T_i,\qquad
 T_i=M_{a^{-1/2}}d_iL_0^{-1}.
\end{equation}
Extend $L_0^{-1}$ by zero on constants. Lemma~\ref{lem:finite} and \eqref{eq:weak}
bound each $\nct(s,T_i^*T_i)$ by

\[
 \Gamma_p(3/4)^Ds^{-p}\sum_xa(x)^{-p}.
\]

For nonnegative matrices, the count of their sum above $t$ is no larger than the
sum of their counts above $t/D$: intersect the complements of their high spectral
subspaces and apply min--max. No commutation is required. With $t=1/c$ in
\eqref{eq:inverse}, inversion on $H_0$ and the single scalar constant mode yield
\begin{equation}\label{eq:count}
 \#\{\lambda(L_a)<c\}\leq
 1+\Gamma_p(3/4)^DD^{p+1}c^pM n^D/m^p.
\end{equation}
No spectral-band dimension is assumed in this derivation.

\begin{proof}[Completion of Theorem~\ref{thm:critical}]
Choose an actual real counting-orthonormal eigenbasis of the real symmetric matrix
$L_a$. Retain the coordinates with eigenvalue $<c$, separately in each of the $h$
real field coordinates. Complexification used in the Fourier proof preserves real
eigenvalue multiplicities; it does not double them. Energy controls coefficients at
eigenvalues $\geq c$, while $c\|C(U)\|^2$ controls the rest. Equation~\eqref{eq:count},
$p=D/2$, and $(1+m)/m\leq2$ give \eqref{eq:rank}. For $n=1$ take the identity on all
$h$ coordinates. This constructs the actual global spectral observation.
\end{proof}

\section{Every lower positive inverse exponent is insufficient}
\begin{theorem}\label{thm:subcritical}
For every fixed $0<r<D/2$, the analogue of Theorem~\ref{thm:critical} under
$\operatorname{mean}((m/a)^r)\leq2$ is false uniformly in the actual scalar class
$1\leq a\leq2m$, $\operatorname{mean}(a)=m$.
\end{theorem}
\begin{proof}
Choose a fixed integer $\ell>3D/c$, and put $P=2\ell+2$, $B=(2\ell+1)^D$.
The triangular window $t(s)=(1-|s|/\ell)_+$ has finite sums
\[
 T=\sum t^2=(2\ell^2+1)/(3\ell),\quad
 V=\sum(d^-t)^2=2/\ell,\quad W=\sum(d^+d^-t)^2=6/\ell^2.
\]
For a product window $\psi_y(x)=\prod_it(x_i-y_i)$ and nonzero actual adjoint $H$,
the actual curl $U_1=d^-_2\psi_y H$, $U_2=-d^-_1\psi_y H$, other components zero,
has backward divergence zero. Its field support and all nonzero forward-gradient
start sites lie in the box $y+[-\ell,\ell]^D$. Direct separable sums give
\[
 \|U\|^2=2VT^{D-1}\|H\|^2,\qquad
 \frac{\E_0(U)}{\|U\|^2}=\frac WV+(D-1)\frac VT
 =\frac3\ell+\frac{6(D-1)}{2\ell^2+1}\leq\frac{3D}\ell<c.
\]
Take sufficiently large multiples $n$ of $P$, set $m=n^2$ and
$J=\lfloor n^{D-2r}/B\rfloor$. These $J$ boxes fit on the separated $P$-grid
eventually, since $P^D\leq Bn^{2r}$. They satisfy
$J\geq n^{D-2r}/(2B)\to\infty$. Put $s=BJ/n^D\leq m^{-r}$ and set $a=1$ on
their union, $a=A=(m-s)/(1-s)$ elsewhere. Eventually $s\leq1/2$, so
\[
 \operatorname{mean}(a)=m,\quad1\leq a\leq2m,\quad A\geq m,
 \qquad\operatorname{mean}((m/a)^r)\leq sm^r+1\leq2.
\]
Each curl's gradients start only in its weak box. The actual $J$-dimensional span
has that same Rayleigh quotient, because both field and gradient supports are
disjoint. Section~1 forces $\dim F\geq J$, whereas the target maximum at $m=n^2$
is $1$. This proves the refutation using original nonlinear data.
\end{proof}
Together the two theorems give the sharp inverse-moment exponent for this specified
family of global scalar conditions. They do not identify all sufficient coefficient
criteria, prove a physical coefficient hypothesis, or impose spatial locality on $C$.

\section{Spatial means and the critical dyadic obstruction}
For a partition $\mathcal P$ into integer rectangles, its actual normalized means are
\[
 C_{\mathcal P}(U)_{R,j}=\sqrt{|R|}\operatorname{mean}_R U_j,
 \qquad \dim F=h\#\mathcal P.
\]
A full dyadic cube tree splits each selected cube into all $2^D$ equal children and
retains unsplit siblings. The origin of the torus coordinate cube may be translated.
This class excludes arbitrary anisotropic rectangle cuts and compressed observations.

\begin{proposition}\label{prop:dyadic}
There are fixed critical-moment profiles with target maximum $1$ for which every
coercive full dyadic cube-tree mean partition has $\Omega(\log n)$ leaves, even after
choosing a translated origin. The same profiles admit an explicit constant-size
rectangular mean partition satisfying the global inequality for every actual cochain.
\end{proposition}
\begin{proof}
Keep $\ell>3D/c$ fixed. Choose a dyadic $L=2^q\geq4\ell+2$, put
$\gamma=\max(1,c/4)$, $n=2^J$ with $J\geq q+1$, $m=\gamma n^2$, and
$S=\{0,\ldots,L-1\}^D$. Set $a=1$ on $S$ and $a=A=(m-s)/(1-s)$ elsewhere,
$s=L^D/n^D$. Then $1\leq a\leq2m$, its mean is exactly $m$, and
\[
 \operatorname{mean}((m/a)^{D/2})\leq1+L^D\gamma^{D/2}=:M_0.
\]
All these parameters except $n$ are fixed, and the target maximum is $1$.

In the hypothetical complete level $j_*=J-q$ of any translated dyadic grid, cells
have side $L$. Each coordinate interval of $S$ meets at most two such cyclic grid
intervals, so one intersection has integer length at least $L/2\geq2\ell+1$.
Taking these intersections independently places a whole gradient-start box for a
compact curl inside $S$ and inside one level-$j_*$ cell. If the proposed partition
had maximum depth at most $j_*$, that cell would lie in one actual leaf. The curl
then has zero means on every leaf, since each component has total sum zero, yet its
weighted Rayleigh quotient is $<c$. This contradicts coercivity. A single zero-extended
window needs no packing divisibility assumption on $n$; $n\geq2L$ suffices here.
The field may depend on the partition, since coercivity is asserted for every field.
Thus maximum depth is at least $J-q+1$, and the full-arity tree identity gives
\begin{equation}\label{eq:dyadiclower}
 \#\mathcal P\geq1+(2^D-1)(J-q+1).
\end{equation}

For a successful unrestricted partition, isolate all $L^D$ vertices of $S$ as
singletons, and use the $D$ disjoint complement rectangles
\[
 R_i=[0,L-1]^{i-1}\times[L,n-1]\times[0,n-1]^{D-i}.
\]
Coverage follows from the unique first coordinate outside $[0,L-1]$.
Every internal edge starts at a strong vertex with coefficient $A$.
On a length-$l$ interval, counting variance and path Cauchy--Schwarz give
\[
 \sum\|u-\bar u\|^2\leq\frac{l(l-1)}{4A}\sum A\|du\|^2.
\]
Indeed the path resistance is at most $(l-1)/A$, and a cut at $t$ is crossed by
$t(l-t)\leq l^2/4$ pairs in the variance identity.
Product-variance tensorization on original fibres yields a rectangle constant
$n^2/(4A)$, with no extra factor $D$. Singletons have zero residual, internal edge
sets are disjoint, and omitted edges have positive energy. Since
$cn^2/(4A)\leq c/(4\gamma)\leq1$, exact orthogonal decomposition proves
\[
 c\|U\|^2\leq Q_w(U)+c\|C_{\mathcal P}(U)\|^2,
 \qquad \#\mathcal P=L^D+D.
\]
These rectangles have potentially unbounded aspect ratio; they leave the cube-tree class.
\end{proof}
The general critical moment question for unrestricted rectangular means remains open.
Changing only the stopping rule within full cube trees cannot resolve it. Classical
weighted Sobolev inequalities \cite{ads} also yield block observations for inverse
moments of finite order strictly greater than $D/2$; the dossier contains the explicit
balanced-box proof. Proposition~\ref{prop:dyadic} explains why its critical tree-counting
loss cannot be removed merely by choosing a better cube-tree stopping criterion.

\section{What a spatial observation proof must still establish}
\label{sec:spatial-obstructions}
The preceding positive spectral result does not construct a rectangular mean map.
Three further checks prevent substitutions for that missing construction. They concern
local geometry, the penalty coefficient, and support of observations, respectively.
They are collected here as limitations of the same dossier, not as additional papers.

\subsection{Critical inverse mass does not remove rectangle geometry}
For a rectangle $R$ of side lengths $l_i$, set $V=\prod_i l_i$, $L=\max_i l_i$,
$\kappa=L^2/V^{2/D}$ and $\mu=\sum_R a^{-D/2}$. The endpoint Sobolev estimate
on balanced cubes, with $q=2D/(D+2)$, transfers to arbitrary rectangles by repeated
coordinate reflection and gives
\begin{equation}\label{eq:aspect}
 \sum_R|u-\bar u_R|^2\leq A_D D^{2/D}\kappa\mu^{2/D}\E_{a,R}(u).
\end{equation}
Here the energy uses original internal edges and their actual starting-vertex weights.
The full reflection multiplicities and graph adapters to the classical estimate
\cite[Section 2.2]{ads} are preserved in the companion geometric note; this is an
analytic adapter, not a new Lean theorem.

The factor cannot be discarded even for uniform coefficients. With $u(x)=x_i-(L-1)/2$
along a longest side, direct sums give
\[
 \sum_Ru^2=V(L^2-1)/12,\qquad
 \E_{a,R}(u)=aV(L-1)/L.
\]
Their ratio is $L(L+1)/(12a)\geq\kappa\mu^{2/D}/12$.
On the thin rectangle of lengths $(L,1,\ldots,1)$, choosing $a=L^{2/D}$ makes
$\mu=1$ while this ratio diverges. Thus compressing a cube-tree chain into thin
rectangles cannot simply reuse a shape-independent inverse-mass certificate.
This prescribed-block test does not defeat every global partition: omitted crossing
edges and a different partition may help. The successful partition in
Proposition~\ref{prop:dyadic} is not contradicted.

\subsection{Fixed and free penalties are different obligations}
Let $C$ be normalized means on a fixed partition of a finite real Euclidean space.
Then $CC^*=I$; $P=C^*C$ is the orthogonal block-mean projection and $Q=I-P$.
For the actual positive stiffness $L_a$, write the blocks on
$\operatorname{ran}P\oplus\operatorname{ran}Q$ as
\[
 A=PL_aP,\qquad B=PL_aQ,\qquad D_0=QL_aQ.
\]
Positivity implies $\ker A\subseteq\ker B^*$: otherwise the energy of $tp+v$
would have a nonzero linear coefficient with zero quadratic coefficient in $t$.
Completing the square with the finite spectral pseudoinverse $A^\dagger$ shows
\begin{equation}\label{eq:shorted}
 L_a+cP\geq cI\quad\Longleftrightarrow\quad
 D_0-B^*A^\dagger B\geq cI\ \hbox{on }\operatorname{ran}Q.
\end{equation}
Indeed the first inequality is exactly $L_a-cQ\geq0$, and its quadratic form is
$\langle p+A^\dagger Bv,A(p+A^\dagger Bv)\rangle+
\langle v,(D_0-cI-B^*A^\dagger B)v\rangle$.
This is classical shorted-operator algebra \cite{anderson}; it supplies no spatial
partition or target count. In particular $D_0\geq cI$ alone is insufficient.

\begin{proposition}\label{prop:formal-cycle}
On the four-cycle with edge weights $(1,6,4,1)$ and blocks $\{0,1\},\{2,3\}$,
coercivity holds on zero block sums with $c=19/5$, but $L_a+cP\geq cI$ fails.
\end{proposition}
\begin{proof}
For $v=(x,-x,y,-y)$, with the counting Euclidean norm,
\[
 \E_a(v)-\frac{19}{5}\|v\|^2
 =\frac{17}{5}\left(x+\frac{35}{17}y\right)^2+\frac{84}{85}y^2\geq0.
\]
The actual normalized observation is
\[
 Cu=\bigl((u_0+u_1)/\sqrt2,\ (u_2+u_3)/\sqrt2\bigr).
\]
For $u=(19,-9,-5,-5)/14$, direct exact evaluation gives
\[
 \E_a(u)=52/7,\quad\|u\|^2=123/49,\quad\|Cu\|^2=25/49,
\]
and hence $\E_a(u)+(19/5)\|Cu\|^2-(19/5)\|u\|^2=-6/35$.
The weights have mean $3$ and lie between $1$ and $6$. Ten focal Lean theorems
preserve the constructed weights, normalized observations, projection identities,
kernel square identity and exact negative witness; Section~\ref{sec:evidence}
identifies the successful immutable source. This is a finite four-cycle result,
not a proved embedding into the full physical divergence-zero kernel.
\end{proof}

For the original class allowing an arbitrary finite $K$, put $J=D_0-cI$.
The necessary and sufficient condition is instead
\begin{equation}\label{eq:free-penalty}
 \exists K\geq0:\ L_a+KP\geq cI
 \quad\Longleftrightarrow\quad J\geq0\ \hbox{and }B\ker J=0.
\end{equation}
Necessity follows by testing $v\in\ker C$, then $v+tp$ with $v\in\ker J$.
For sufficiency, $B\ker J=0$ implies $B^*p\in\operatorname{ran}J$.
Completing the square in $v+J^\dagger B^*p$ leaves
$A+(K-c)I-BJ^\dagger B^*$, so $K=c+\|BJ^\dagger B^*\|$ suffices.
In particular a proved all-cochain margin $D_0\geq(c+\epsilon)I$ gives
$K=c+\|L_a\|^2/\epsilon$. For actual $a\leq2m$, the elementary edge estimate
$|r-s|^2\leq2|r|^2+2|s|^2$ gives $\|L_a\|\leq8Dm$, including the parallel
edge occurrences at $n=2$. A margin $2c$ therefore permits
$K=c+64D^2m^2/c$, without increasing the number of observations.
No margin or partition is proved by this conditional implication. Ordinary block
averaging need not preserve divergence, so an estimate only on the physical kernel
cannot silently replace the all-cochain hypothesis in this decomposition.

Even free $K$ needs care at equality. On the same cycle let
$s=\sqrt{85}$, $c_*=(17-s)/2$ and $v=(7,-7,6-s,s-6)$. Its kernel defect at $c_*$
is zero, since for general $(x,-x,y,-y)$ it equals
$(s-6)[x+7y/(s-6)]^2$. For
$p=(\delta/2,\delta/2,-\delta/2,-\delta/2)$ the full defect is
\[
 \E_a(v+p)+K\|C(v+p)\|^2-c_*\|v+p\|^2
 =(7+K-c_*)\delta^2-10(13-s)\delta.
\]
For each finite $K\geq0$ a sufficiently small positive $\delta$ makes it negative.
This additional equality example is exact analytic algebra and finite diagnostic
evidence, not one of the ten Lean theorems of Proposition~\ref{prop:formal-cycle}.

\subsection{Sparse point samples cannot replace stable averages}
For uniform $a=m$, if an observation depends only on vertex values on a set $S$, $|S|=s$, and
$C(0)=0$, then the scalar field $u=1$ off $S$ and $0$ on $S$ has zero observation.
At most $2Ds$ forward-edge occurrences cross $S$, so
$\E_m(u)\leq2Dm s$. Consequently all-cochain coercivity requires
\[
 s\geq\frac{cn^D}{2Dm+c}.
\]
For an actual divergence-zero test, fix a component direction $j$ and let the field
be constant along every $j$-line, equal to zero on each line meeting $S$ and one on
the other lines, in one fixed internal coordinate. There are at most $s$ hit lines.
Its backward divergence is zero, its norm squared is at least
$n(n^{D-1}-s)$, and transverse crossing energy is at most $2(D-1)mns$.
Thus necessarily
\[
 s\geq\frac{cn^{D-1}}{2(D-1)m+c}.
\]
If every transverse line is hit, the second bound is immediate without a nonzero test. These bounds give the required obstruction:
at $a=m=n^2$ and critical moment $M=1$, bounded samples fail for all cochains
when $D\geq3$ and for these physical tests when $D\geq4$.
For the same uniform coefficients one global mean succeeds for $0<c\leq16$,
by the nonzero Fourier gap $n^2\lambda(k)\geq16$.
This comparison rules out a point-sampling shortcut, not rectangular averages or
singletons supplemented by complementary block means. A nonlinear local observation
is tested on sufficiently small multiples of these fields.

\section{Evidence provenance and verification limits}\label{sec:evidence}
The analytic sources, their exact revisions and the executable finite diagnostics are
collected in a hash-allowlisted companion package. Its map gives full SHAs and SHA256
values; short forms below are identifiers, not substitute hashes.
\begin{center}\small
\begin{tabular}{lll}
Item & Source revision & Verification status\\\hline
Actual flat Hodge identity & \texttt{dfd49d0b0730} & earlier formal source\\
Critical spectral theorem & \texttt{527fcdcd119a} & analytic adapters\\
Dyadic spatial separation & \texttt{fb65a67af360} & analytic adapters\\
Higher moment block proof & \texttt{6dcb893155e5} & analytic adapters\\
Exact finite fixture code & \texttt{4bb460e1537e} & own Colab109\\
Full diagnostic report & \texttt{602ae3f0b59a} & finite rational batch
\end{tabular}
\end{center}
Colab109 checked four actual profiles on $(\mathbb Z/2\mathbb Z)^4$ using rational
arithmetic. Each normal and optimized Python mode accepted four profiles and rejected
16 real matrix, inverse or flow mutations, with matching outputs and exit code zero.
It checked the constant mode, forward symbols, actual pseudoinverse identities,
feasible/minimizing flow energy, comparison inertia, strict low-spectrum multiplicities
and noncommuting counts. There are 12 distinct density examples, evaluated again for
each profile: 48 evaluations per mode, not 48 independent examples.

All eight files in the diagnostic archive's inner manifest were rehashed. Runtime
execution was recorded at 17:13:06.300923--17:13:16.201862 UTC on 2 November2026,
with automatic runtime release and zero sessions verified afterwards. Rounded resource
units do not prove zero cost. No local Windows mathematical execution occurred.
This single embedded-SHA finite batch is not two fresh matching clones, a Lean proof,
a theorem about arbitrary dimensions or odd lattices, or independent mathematical review.
No new Lean run occurs in this appendix's PDF assembly. Operational repository checks
on unchanged Lean files do not verify these new analytic proofs.

There is also subsequent, strictly focal formal evidence. Source
\texttt{41f8492c6190} supplies the ten four-cycle theorems and their oracle.
Own Colab118 compiled that same source in two fresh clones; both actual oracle logs
and compiled module hashes agree. Each of four normal/optimized validator runs
accepted ten readouts and rejected three real corruptions. All readouts use only
\texttt{propext}, \texttt{Classical.choice}, and \texttt{Quot.sound}.
The full archive has 69 manifest files and is included with the failed Colab116/117
compiler attempts. Source hashes, complete logs and exact reports are allowlisted.
This verifies Proposition~\ref{prop:formal-cycle}'s finite identities, not the
geometric, shorted-operator, free-penalty or general rectangular adapters.
Own Colab115 separately checked 18 finite diagnostics and eight real mutations in
each mode of two fresh clones. Finite examples do not establish their general claims.

\section{Model assessments and review context}
The prior analytic-source first self-assessments were 6/10 for the critical spectral
note, 5/10 for the dyadic note, and 5/10 for the exact finite diagnostic package.
They apply only to their own exact source versions. All were by the preparing Codex
agent; the served backend identity and effort were unexposed. They are neither independent
model reviews nor human peer review. No earlier score is inherited by this appendix.
The previous appendix PDF and its first exact assessment remain historical artifacts;
their assessment does not transfer to this consolidated revision. This revision's
first exact-PDF review is a separate SHA-linked sidecar after visual inspection.
No grade-seeking repeated review,
editorial acceptance or publication status is implied.

The principal limitations and objections are these. The new finite Fourier, inverse,
field and partition adapters are analytic rather than kernel-verified. The classical
Rumin--Frank method and elementary refinement geometry limit any novelty claim; a
complete priority comparison is unfinished. The sufficient theorem uses one shared
scalar coefficient profile, and its distribution condition is not proved from arithmetic
means. Spectral coordinates are global; unrestricted critical rectangular means remain
unresolved. Exact finite diagnostics cover a small selected family, and some count bounds
are loose. The constants are safe rather than sharp. This appendix makes no statement
about an interacting or continuum physical mass gap. PR29 remains 20/41 and TermSource=0.

\section{The single remaining bounded question}
For fixed $D\geq3,c>0,N_c\geq2,M$, do all actual scalar profiles
$1\leq a\leq2m$, exact mean $m$, satisfying \eqref{eq:moment}, admit
\eqref{eq:rank} using normalized means on an unrestricted axis-aligned integer
rectangular partition? The mechanism must leave full cube trees, for example through
anisotropic cuts with a proved variance bound. A falsifier must defeat every such
rectangular partition, not merely every dyadic cube tree. After three substantive units
without a proved reduction or useful falsifier, the route must be changed or its concrete
blocker identified. Formalization and consolidation of the existing results are preserved
work, not additional active new research questions.

\begin{thebibliography}{9}
\bibitem{anderson} W.~N. Anderson, Jr. and G.~E. Trapp,
\emph{Shorted Operators. II}, SIAM J. Appl. Math. 28(1) (1975), 60--71.
\url{https://doi.org/10.1137/0128007}.
\bibitem{frank} R.~L. Frank, \emph{Cwikel's theorem and the CLR inequality},
arXiv:1206.3325v1 (2012), Section2, Theorem2.1, Corollary2.2, Lemmas2.4--2.5.
\url{https://arxiv.org/abs/1206.3325}.
\bibitem{ads} S.~Andres, J.-D. Deuschel and M.~Slowik,
\emph{Harnack inequalities on weighted graphs and some applications to the random
conductance model}, arXiv:1312.5473v6 (2019), Section2.2, equation(2.5) and Proposition2.1.
\url{https://arxiv.org/abs/1312.5473}.
\end{thebibliography}
\end{document}
